\scanpage{96}

\deleted{On a, si $N$ \unsure{$<$|$\leq$}}

\[
\deleted{\ell^1\subset\ell^N\subset\ell^\infty}\tag*{\deleted{(6)}}
\]

\deleted{et la norme de $\ell^1$ est plus fine que celle induite par $\ell^N$, \unsure{qui est} et la norme de $\ell^N$ est moins fine que celle induite par}

Si $N_{\deleted{\unsure{1}}}'$ et $N_{\deleted{\unsure{2}}}''$ sont deux normes de Schatten telles que $N_{\deleted{1}}'\leq N_{\deleted{2}}''$, il est immédiat qu'on a $\ell^{N'_{\deleted{1}}}\subset\ell^{N''_{\deleted{2}}}$, et la norme de $\ell^{N'}$ \deleted{$\unsure{\ell^{N_1}\subset\ell^{N_2}|\ell^{N'}\subset\ell^{N''}}$} est plus fine que celle induite par $\ell^{N''}$.

\deleted{En particulier,} En particulier, si on remarque que toute norme de Schatten est comprise entre les normes de Schatten $N_1(x)=\sum_i |x_i|$ \deleted{\unsure{et}} $N_\infty(x)=\underset{i}{\operatorname{Sup}}\,|x_i|$, et que les espaces $\ell^N$ associés à ces dernières sont \deleted{$\unsure{\ell^1|\ell^{N_1}}$ et $\unsure{\ell^N|\ell^{N_\infty}}$} $\ell^1$ et $\ell^\infty$, \deleted{\unsure{est on}} voit que

\[
\ell^1\subset\ell^N\subset\ell^\infty\tag*{(6)}
\]

les applications identiques \deleted{\unsure{de $\ell^1$}} $\ell^1\to\ell^N$ et $\ell^N\to\ell^\infty$ étant de norme $\leq 1$.

\deleted{Notons la formule}

\[
\deleted{\sum\quad N\left(\sum_n x_{\unsure{i}}^{k}\right)\leq\sum_k N(x^k)}\tag*{\deleted{(7)}}
\]

\deleted{valable pour toute suite $\unsure{x^i=(x_i^i)|x^k}$ de suites $x^k\in C^N$, qui soit \deleted{\unsure{convergente}} sommable pour chaque coordonnée $i$\,: $\sum_k|x_i^k|<+\infty$. Il suffit de prouver (7) si le dernier mb est fini\,: alors, en posant $x=\sum_k x^k$, on aura $|x|_m\leq\sum_k|x^k|_m$ d'où $N(|x|_m)\leq N(\sum_k|x^k|_m)$}

% slanted
\margin{\unsure{Si Si}}

\deleted{or \inserted{$N$ étant une fonction continue sur $R^m$, on aura} $N(\sum_k|x^k|_m)=\lim_{\unsure{k|n}\to\infty}N(\sum_{k=1}^n|x^k|_m)$}

\deleted{et comme $N(\sum_{k=1}^{n}|x^k|_m)\leq\sum_{k=1}^n N(|x^k|_m)\leq\sum_{k=1}^{\unsure{\infty|n}} N(x^k)$, on aura $N(\sum_k|x^k|_m)\leq\sum_k N(x^k)$, donc pour tout $m$ on a $N(|x|_m)\leq\sum_k N(x^k)$, d'où par définition de $N(x)$ l'inégalité (7).}
% archive sheet number: 97
